\scanpage{49}
% Pencil archival foliation: 48.
provienne d'un $K$-hom
$R\unsure{\hookrightarrow|\longrightarrow}R'[s_1,\ldots,s_N]$
($R,R'$ ext. de type fini de $K$),
que l'isomorphie du système des
$(\gamma_i(R))$ et du système des $(\gamma_i(R'))$
implique que $R$ et $R'$ sont ``équivalents''
en ce sens que chacun s'immerse
dans une extension transcendante pure
de l'autre ----

\uline{Exemple} 1) Soient $C,C'$ deux courbes
\inserted{lisses complètes \unsure{conn.}}
\inserted{sur $k$ alg. clos,}
supposons que les jacobiennes soient telles
qu'il existe un hom $J\to J'$ à
noyau fini. Prouver qu'il existe
\deleted{C'} une application rationnelle dominante
\[
C'[t_1,\ldots,t_N]\unsure{-|\longrightarrow}C\quad ??
\]
\margin{\uline{$k$ alg. clos} \uline{au besoin!}}

2) Si les $\underline{\unsure{\gamma}^{\,i}}(R)$ sont nuls pour
$i\unsure{>|\geq}0$,
prouver que $R$ est \unsure{unirationnelle|unirationnel} ??
[En particulier, si la cohomologie
\inserted{de $X$ var. proj. non sing.}
entièrement \unsure{sphérique} ??]
\margin{Cas d'une surface ?? [\unsure{coh. alg.!}]}
